%% ============================================================
%%  §1  Introduction
%% ============================================================

\section{Introduction}
\label{sec:intro}

Modern software stacks mix languages. A backend service might be written in
Rust for performance-critical subsystems, TypeScript for the API layer, and
Python for tooling. When the goal is formal verification, this heterogeneity
poses a quiet but persistent problem: each existing verification framework is
coupled to one language. Aeneas~\cite{ho2022aeneas} targets safe Rust through
functional translation. Verus~\cite{verus2023} provides deductive verification
for Rust via SMT-backed contracts. Creusot~\cite{creusot2022} translates Rust
into Why3 proof obligations. None of these tools say anything about TypeScript,
and none of their specifications are portable across language boundaries.

The result is that proving a property about a Rust function and then extending
that proof to a semantically equivalent TypeScript function requires starting
over: a different tool, a different specification language, a different proof
style. Cross-language correctness — the claim that two implementations of the
same function, in different languages, satisfy the same functional
specification — is effectively unprovable with the current toolkit.

\subsection{The Key Insight}

Rust and TypeScript differ structurally. Rust's ownership system enforces
linear resource discipline at compile time; TypeScript runs on a garbage-collected
runtime with dynamic typing and implicit coercions. Their compilers produce
radically different intermediate representations. Yet at the level of pure
functional semantics, both execute sequences of operations that either
succeed (producing a value) or fail (raising an exception or panicking). That
shared structure is precisely what a state-error monad captures.

If programs from both languages are translated into a single monadic target —
a monad that models local state and runtime failure — then a specification
written over that monad describes both programs without modification. This
observation drives the entire design of \emph{lean-pl-verify}.

\subsection{Our Approach}

We define two components that are intentionally language-agnostic:

\begin{itemize}[leftmargin=1.5em]
  \item \textbf{\RustM}, the execution monad:
        $\RustM\;\sigma\;\alpha \;=\; \texttt{StateT}\;\sigma\;(\texttt{Except}\;\texttt{PanicReason})\;\alpha$.
        This represents a computation that may read and update a state $\sigma$,
        may return a value of type $\alpha$, and may instead signal a panic.

  \item \textbf{\ProgramSpec}, the specification language: an inductive type with
        five constructors (\texttt{pureOutput}, \texttt{nocrash}, \texttt{both},
        \texttt{withFuel}, \texttt{precond}) and a satisfaction relation
        $m \satisfies P$ defined by recursion on $P$.
\end{itemize}

Programs from Rust and TypeScript are translated into \RustM via separate deep
embeddings. A \emph{deep embedding} means that the source language's syntax is
represented as an inductive Lean~4 datatype, and execution is defined as a
definitional function over that datatype. The Lean kernel can therefore evaluate
closed programs directly, which is what enables \rfl-style proofs for pure
functions. No external trusted component sits between the source program and
the Lean kernel.

For Rust we embed a substantial fragment of LLBC (Low-Level Borrow
Calculus)~\cite{charon2024}: the AST covers assignments, conditionals,
sequential composition, loops with \texttt{break}/\texttt{continue},
and explicit function calls. For TypeScript we embed a parallel AST covering
expression evaluation, binary operations, conditionals, and function calls.
Both ASTs use the same \Locals representation (a list of \texttt{Value}s),
the same \RustM monad, and the same \ProgramSpec specification language.

\subsection{Contributions}

This paper makes the following contributions.

\begin{enumerate}[label=\textbf{C\arabic*.},leftmargin=2em]
  \item \textbf{RustM monad} (\S\ref{sec:monad}). A kernel-transparent
        state-error monad with a \texttt{PanicReason} inductive type covering
        arithmetic overflow, division by zero, explicit panics, and index
        errors. Two proved monad laws (\texttt{bind\_ok}, \texttt{pure\_ok})
        and a forward-backward model for mutable references.

  \item \textbf{Deep LLBC embedding} (\S\ref{sec:llbc}). An inductive Lean~4
        encoding of the LLBC AST with a fuel-parameterized interpreter
        \texttt{evalStmtFuel} that is definitionally transparent to the kernel.
        The interpreter handles the full LLBC fragment used by Charon~\cite{charon2024}
        for safe Rust: loops, break, conditionals, multi-function environments.

  \item \textbf{ProgramSpec} (\S\ref{sec:programspec}). A language-agnostic
        specification combinator library with 17 kernel-checked theorems
        (\texttt{sat\_pure}, \texttt{sat\_bind}, \texttt{sat\_both\_intro},
        \texttt{sat\_pureOutput\_mono}, and~13 others). The library makes no
        reference to LLBC or TypeScript ASTs.

  \item \textbf{Forward-backward ownership model} (\S\ref{sec:monad}). A pure
        functional encoding of \texttt{\&mut T} via \texttt{FBPair} with four
        identity laws proved in Lean, matching the formal semantics of
        Aeneas~\cite{ho2022aeneas} without requiring separation logic.

  \item \textbf{Loop invariant methodology} (\S\ref{sec:rust-proofs}). A
        reusable inductive invariant pattern demonstrated on three loop
        functions: \texttt{sumTo} (Gauss formula), \texttt{fact} (factorial),
        and \texttt{fib} (Fibonacci, \S\ref{sec:fib}). The pattern involves an
        explicit overflow predicate, a parametric loop-correct theorem proved
        by induction on remaining iterations, and a safety-implies-overflow
        bridge lemma.

  \item \textbf{Proof automation} (\S\ref{sec:automation}). Four Lean~4 tactic
        macros (\texttt{llbc\_verify}, \texttt{llbc\_verify\_prop},
        \texttt{llbc\_verify\_cond}, \texttt{llbc\_verify\_loop}) that reduce
        typical 12-line LLBC proofs to single-line calls. Measured compression:
        12$\times$.

  \item \textbf{Relational semantics and full adequacy} (\S\ref{sec:semantics}).
        A relational big-step semantics (\texttt{EvalStmt}, \texttt{EvalFun},
        and four sub-relations) provides a language-independent ground truth for
        what LLBC programs mean. A \emph{full adequacy} theorem (18 kernel-checked
        proofs, 0 \sorry) establishes both soundness (every interpreter output is
        relationally derivable) and completeness (every relational derivation is
        witnessed by some fuel bound), closing the gap between the fuel-based
        evaluator and a declarative specification.

  \item \textbf{Charon end-to-end pipeline and real-crate extraction}
        (\S\ref{sec:eval-charon}--\ref{sec:eval-numinteger}).
        A 463-line Python translator (\texttt{charon2lean.py}) automatically
        converts Charon LLBC JSON to \texttt{LLBCFunDef} Lean terms.
        57 theorems prove the Charon-generated definitions correct across
        18 functions, including \texttt{pow} and \texttt{gcd}.
        Additionally, 5 theorems verify \texttt{Integer::is\_even<u32>}
        extracted from \texttt{num-integer} v0.1.45 (${\approx}$45M downloads),
        a published crate on \url{crates.io}.
        The pipeline: \texttt{verified\_fns.rs} $\to$ Charon $\to$ LLBC JSON $\to$
        \texttt{charon2lean.py} $\to$ \texttt{CharonDefs.lean} $\to$ proofs.

  \item \textbf{TypeScript extension and unification} (\S\ref{sec:typescript}).
        33 theorems for a TypeScript AST interpreter (including a verified
        \texttt{while} loop) using the same \RustM monad and \ProgramSpec
        specification. Six cross-language unification theorems — covering
        \texttt{max}, \texttt{min}, \texttt{mul}, \texttt{add}, \texttt{neg},
        and the looping \texttt{sum\_to} — formally assert that Rust and
        TypeScript implementations satisfy identical specifications.

  \item \textbf{Reproducible artifact} (\S\ref{sec:evaluation}). \textbf{258} kernel-checked
        theorems in 5\,361 lines of Lean~4 code, \textbf{0 sorry}. All theorems are verified
        by a single command: \lakeCmd{lake build Theorems}. The artifact is
        independently reproducible given Lean~4.30.0-rc2 and Mathlib4.
\end{enumerate}

\subsection{What This Work Is Not}

Honesty about scope matters. This work does not verify unsafe Rust: pointer
arithmetic, raw memory, and concurrency are outside the model.
The verification covers functional correctness properties (\eg correctness of
Fibonacci), not security properties or liveness.
The Charon pipeline handles the 16 functions in the artifact; scaling to large
real-world crates is future work.
All 258 theorems are fully kernel-checked: \textbf{0 sorry}.

\subsection{Why Deep Embedding?}

The prevailing approach in tools like Aeneas uses a \emph{shallow} embedding:
Rust functions are compiled to native Lean~4 terms, bypassing an explicit
interpreter. This produces cleaner generated code and shorter proofs for pure
functions. The cost is a trusted code generator: correctness of the generated
Lean depends on correctness of the compiler, which is not itself formally
verified.

Our deep embedding inverts this tradeoff. The interpreter is a Lean~4 \texttt{def},
so every computation step is visible to the Lean kernel. No external component
is trusted. The cost is that proof terms are larger, and loop invariant proofs
require more manual effort. We believe this tradeoff is the right one for a
library whose primary purpose is to establish the \emph{foundations} of
cross-language verification: when the question is ``what do these programs
actually mean?'', having the answer be ``exactly what the Lean kernel says''
is a valuable guarantee.

\subsection{Paper Organisation}

Section~\ref{sec:background} reviews background and related work.
Section~\ref{sec:monad} introduces \RustM and the forward-backward model.
Section~\ref{sec:llbc} presents the LLBC embedding and interpreter.
Section~\ref{sec:programspec} defines \ProgramSpec.
Section~\ref{sec:rust-proofs} walks through the Rust verification results,
including the loop invariant methodology.
Section~\ref{sec:automation} describes the proof automation macros.
Section~\ref{sec:typescript} covers the TypeScript extension and
cross-language unification.
Section~\ref{sec:fib} presents the Fibonacci case study.
Section~\ref{sec:semantics} presents the relational semantics and full adequacy proof.
Section~\ref{sec:evaluation} gives the evaluation.
Section~\ref{sec:agentic} reflects on the agentic development process.
Section~\ref{sec:related} discusses related work.
Section~\ref{sec:conclusion} concludes.
